\documentclass[10pt,a4paper]{amsart}
\usepackage[margin=19mm]{geometry}
\usepackage{amsmath,amssymb,mathtools}
\usepackage[scr=rsfs,cal=euler]{mathalfa}
\usepackage{xfrac}
\usepackage[unicode,hidelinks,bookmarks=false]{hyperref}
\newcommand{\ie}{i.e.\ }
\newcommand{\cf}{cf.\ }
\newcommand{\resp}{respectively\ }
\newcommand{\Subsection}{\subsection}
\providecommand{\nobreakdash}{\nobreak\hbox{-}}
\setlength{\emergencystretch}{2em}
\multlinegap0pt
\usepackage{amssymb}
\newtheorem{prop}{Proposition}
\newtheorem{claim}{Claim}
\newcommand{\HF}{\textnormal{HF}}
\newcommand{\inv}{^{-1}}
\newcommand{\res}{\upharpoonright}
\title{Computable vs Descriptive Combinatorics of Local Problems on Trees}
\author{Felix Weilacher}
\date{Source: arXiv:2208.06689v2 (CC BY 4.0)}
\begin{document}
\maketitle

\noindent\emph{Source and licence.} Computable vs Descriptive Combinatorics of Local Problems on Trees. Version arXiv:2208.06689v2 (CC BY 4.0), distributed by arXiv under the Creative Commons Attribution 4.0 International licence, http://creativecommons.org/licenses/by/4.0/, as recorded in the arXiv licence metadata for this exact version and shown on the abstract page https://arxiv.org/abs/2208.06689v2. Full licence notice: \texttt{licenses/arxiv-2208.06689-CC-BY-4.0.txt}.\par\smallskip

\noindent\textit{Editorial scope.} Counted unit: the article's prop prop:hcconstruction (source0.tex lines 173--175). The article's complete original proof of it is delivered with it (source0.tex lines 177--264, one \texttt{proof} environment of 88 lines). No mathematics is written, repaired or re-derived: the statement and the proof are the article's own lines, in the article's own notation, and no retained line is substituted or re-flowed.  No further source-local context is needed: every cross-reference of every retained line resolves inside the two delivered spans.  One declared adaptation: exactly 3 ancillary strings inside the retained spans are omitted (image-inclusion commands and, where the source repeats a label, the repeated label definitions) and no other line differs from the source; the figure environments, with their placement, captions and labels, are kept, and the package records the omitted strings in fidelity receipts bound to the affected bodies.  No retained line contains a citation, so the wrapper carries no bibliography and no bibliography entry.  The retained lines contain the article's own diagram code, which the wrapper typesets with the standard tikz packages; no image file is delivered and the package declares no asset dependency.  Editorial formatting only, in addition: an AMS \texttt{amsart} preamble that restates the article's own declarations and macros used by the retained lines, namely the environments \texttt{prop}, \texttt{claim} on separate counters, together with the standard packages those lines need.  The retained environments print in this document as Proposition 1, Claim 1 in order of appearance, whereas the article prints the same items with the numbering of its own full document.  The complete native source of the article as distributed in the arXiv e-print for this exact version is bundled unchanged as \texttt{sources/129440/source0.tex}.
\begin{prop}\label{prop:hcconstruction}
        If an LCL on $\Delta$-regular graphs $\Pi = (\Sigma,\mathcal{V},\mathcal{E})$ is not full, there is a computable truly $\Delta$-regular forest (thus highly computable) with no computable $\Pi$-coloring.
\end{prop}

\begin{proof}

If $\Pi$ is not full, then for every nonempty $\mathcal{V}' \subset \mathcal{V}$ and $l' \geq 1$, there are $l \geq l'$, multisets, say $a,b \in \mathcal{V}'$, and labels, say $\alpha, \beta \in \Sigma$ appearing in $a$ and $b$ respectively such that the following holds: Let $P$ be a $\Delta$-regular path of length $l$, say with vertices $x_0,\ldots,x_l$ in order. Label the half edges incident to $x_0$ such that $c(x_0) = a$ and $c(x_0,x_1) = \alpha$, and likewise for $x_l, b, x_{l-1}$ and $\beta$ respectively. Then if this labeling is extended to a $\Pi$-coloring $c$ of $P$, some vertex $x_i$ has $c(x_i) \not\in \mathcal{V}'$. Note that if $\mathcal{V}' = \mathcal{V}$, this means there is no such extension. 

Since $\Sigma$ is finite, we can say instead that for all nonempty $\mathcal{V}' \subset \mathcal{V}$, there exists $a,b \in \mathcal{V}'$ and $\alpha,\beta \in \Sigma$ which have this property for infinitely many $l$. Fix choices of such $a,b,\alpha$, and $\beta$ for each $\mathcal{V}'$, and call them \textit{bad} for $\mathcal{V}'$. Also call the infinite set of $l$ witnessing this badness the set of \textit{bad} $l$ for $\mathcal{V}'$. Note that determining whether a given $l$ is bad for a given $\mathcal{V}'$ is computable: It can be done by checking all possible colorings of a $\Delta$-regular path of length $l$ with labels from $\Sigma$, of which there are only finitely many.

%%%%%%%%%%%%%%%%%%%%%%%%%%%%%%%%%%%%%

Fix $X,E \subset \HF$ disjoint infinite computable sets, and a computable well order of $\HF$ of type $\omega$. $X$ and $E$ will be our vertex and edge sets respectively. Fix an effective enumeration $\{\phi_t \mid t \in \omega \}$ of the set of partial recursive functions $X \times E \rightarrow \Sigma$. (These are candidates for computable colorings of the half edges of our graph.) Fix also an effective enumeration $\{t_n \mid n \in \omega \}$ of $\omega$ which lists each natural number infinitely many times. We are about to describe a recursive construction of $\Delta$-regular forests $\emptyset = T_0 \subset T_1 \subset \cdots$ with vertex sets contained in $X$. The vertex set for each $T_n$ will be finite. It will happen that in our construction every $x \in X$ will eventually become a vertex and be made to have degree $\Delta$. The limit $T = \bigcup_n T_n$ will then be a computable truly $\Delta$-regular forest. Our construction will be designed so that $T$ does not admit a computable $\Pi$-coloring.

Alongside $T$, we will recursively build a $T$-invariant total computable function $C:X \rightarrow \omega$. The idea is that the vertices in $C\inv(t)$ will be responsible for ensuring that $\phi_t$ is not a $\Pi$-coloring of $T$. At stage $n$, the vertices on which $C$ has been defined will be exactly those in the vertex set of $T_n$. 

At a typical stage in the construction, we will take currently unused (not in the vertex set of our current finite tree) vertices, define $C$ on them, and give them $\Delta$ incident half edges. We will refer to these as ``new vertices'', and always choose the least (according to our previously fixed order of $\HF$) available. For example, if $x$ has degree 0 in $T_n$, and in the definition of $T_{n+1}$ we say ``add $\Delta$ new vertices as neighbors to $x$'', it means: let $y_0,\ldots,y_{\Delta-1}$ be the least elements of $X$ not in the vertex set of $T_n$, let $e_0,\ldots,e_{\Delta-1}$ be the virtual edges incident to $x$ in $T_n$, and add the half edges $(y_i,e_i)$ to $T_{n+1}$ for $i < \Delta-1$. The exact same treatment will be given to edges. For example, in the previous example, we would probably want to finish by saying ``add $\Delta-1$ new virtual edges incident to each $y_i$'' to maintain $\Delta$-regularity. This would mean: let $f_{i,j}$ for $i < \Delta$, $j < \Delta - 1$ be the least elements of $E$ not in the edge set of $T_n$, and add the half edges $(y_i,f_{i,j})$ to $T_{n+1}$ for each $i,j$. 

In both steps of the previous example, our language leaves some ambiguity about exactly which new vertices/edges are connected to which old edges/vertices. The exact decision here will never matter (except of course in that it should be made in a computable way).

If the stage of our construction is $n$, $C$ will always be defined to be $t_n$ on new vertices. Thus, our construction will essentially build $\omega$ forests in parallel, one for each $t$. Let us fix $t$, and describe the $\omega$ steps in the construction of the $t$-th forest, i.e, $T \res C\inv(t)$.

\begin{figure}
    \centering
    
    \caption{Stage 0 of the construction for a fixed $t$ for $\Delta = 3$.}
    \label{fig:0}
\end{figure}

Fix a very large $N_0 >> |\Sigma|$ to be determined later, independent of $t$. In the 0-th stage of the construction, place $N_0$ new vertices, say $x_0,\ldots,x_{N_0-1}$, in $C\inv(t)$. Also give each $x_i$ $\Delta$ new virtual incident edges. This is shown in Figure \ref{fig:0}.

We also initialize several variables: Set $\mathcal{V}' = \emptyset$ and $N = N_0$. Also, for each $i < N$, let $P_i$ be the $\Delta$-regular path of length 0 whose unique vertex is $x_i$.

At the start of an arbitrary successor stage, say $m+1$, suppose we have some $\mathcal{V}' \subset \Sigma$ and $N \in \omega$ very large. Suppose also that we have $\Delta$-regular paths $P_i$ for $i < N$ and vertices $y_i^c$ for $i < N$ and $c \in \mathcal{V}'$ such that
\begin{enumerate}
    \item The $P_i$'s and $y_i^c$'s all lie in distinct $C\inv(t)$-components.
    \item For each $i$, if $\phi_t$ restricted to the half edges meeting $P_i$ is a $\Pi$-coloring of $P_i$, there must be some vertex $x$ of $P_i$ with $\phi_t(x) \not\in \mathcal{V}'$.
    \item $\phi_t(y_i^c) = c$ for each $i$ and $c$.
\end{enumerate}
Note that this is all satisfied after stage 0. 

\begin{figure}
    \centering
    
    \caption{The ``uninteresting case'' for a step of the construction for a fixed $t$, for $\Delta = 3$.}
    \label{fig:default}
\end{figure}

We start our successor stage by running the Turing machine computing $\phi_t$ for $m$ steps on all of the half edges meeting all of the $P_i$'s. 

If $\phi_t$ fails to converge in $\leq m$ steps on one of these inputs, or if it converges on all of them but the result fails to be a $\Pi$-coloring of some $P_i$, we do what is shown in Figure \ref{fig:default}: For each virtual edge $e$ with unique endpoint in $C\inv(t)$, we make $e$ true by adding a new vertex, say $x$, as its other endpoint, then to maintain $\Delta$-regularity we add $\Delta-1$ new edges as virtual edges incident to this $x$. We call this the ``uninteresting case''.

The ``interesting case'' is of course if $\phi_t$ does converge on all these inputs in $\leq m$ steps, and the result is a $\Pi$-coloring for each $P_i$. In this case, by condition (2) above, there is a vertex in each $P_i$, say $z_i$, with $\phi_t(z_i) \not\in \mathcal{V}'$. Let $d \in \mathcal{V} \setminus \mathcal{V}'$ such that $\phi_t(z_i) = d$ for at least $N/|\mathcal{V}|$ $i$'s. Update $\mathcal{V}'$ to $\mathcal{V}' \cup \{d\}$.

Now $\mathcal{V}'$ is nonempty, so let $a,b \in \mathcal{V}'$ with $\alpha \in a$ and $\beta \in b$ be bad for $\mathcal{V}'$. Between the $z_i$'s and the $y_i^c$'s, and using condition (1) above, we can choose degree $\Delta$ vertices $v_i$ and $w_i$ for $i < M := N/(2|\mathcal{V}|)$ such that $\phi_t(v_i) = a$ for each $i$, $\phi_t(w_i) = b$ for each $i$, and all $2M$ of these vertices lie in distinct $C\inv(t)$-components. (The factor of $1/2$ is needed for the case $a = b$.)

\begin{figure}
    \centering
    
    \caption{The ``interesting case'' for a step of the construction for a fixed $t$. The top and bottom of the image represent the before and after states respectively. The dashed line encloses one of the new paths, $P_i$.}
    \label{fig:join}
\end{figure}

Now for each $i < M / 2$, pick $e_i$ and $f_i$ virtual edges meeting the components of $v_i$ and $w_i$ respectively such that if $e$ is the first edge in the path from $v_i$ to $e_i$, then $\phi_t(v_i,e) = \alpha$, and likewise for $w_i$, $f_i$, and $\beta$. Note that these exist since our components are finite $\Delta$-regular trees. As shown in in the first step in Figure \ref{fig:join}, join $e_i$ and $f_i$ with a path of new vertices so that the resulting path between $v_i$ and $w_i$ has a total length which is bad for $\mathcal{V}'$. (Recall that this is possible as there are arbitrarily long bad lengths, and that determining if a length is bad is computable.) Call this new path $P_i$. Also, as shown in the second step in Figure \ref{fig:join}, add $\Delta-2$ new virtual incident edges to the newly added vertices in $P_i$ to maintain $\Delta$-regularity. Condition (2) is satisfied for each $P_i$ by definition of bad.

We can now finish this stage of the construction with the updated value of $N$ being $M/2$. We still need to define the new $y_i^c$'s. Use the unused $v_i$'s and $w_i$'s for $c \in \{a,b\}$. For $c \not\in \{a,b\}$, use the first $M/2$ $z_i$'s with $\phi_t(z_i) = d$ if $c = d$, and the first $M/2$ $y_i^{c}$'s from the previous stage of the construction if $c \neq d$ (i.e, if $c$ was in the previous $\mathcal{V}'$.) (1) is clear by construction.

Thus the description of the construction is complete. We can also see now that the correct value of $N_0$ to take was $N_0 > (4|\mathcal{V}|)^{|\mathcal{V}|}$, since the value of $N$ was divided by $4|\mathcal{V}|$ each time the interesting case occurred, and it cannot occur more than $|\mathcal{V}|$ times for a fixed $t$ (Since $|\mathcal{V}'|$ increases by one each time it does).


\begin{claim}\label{claim:computable}
$T$ is a computable truly $\Delta$-regular forest.
\end{claim}
\begin{proof}
At each stage in the construction, we either added new trees as components, added new degree 1 neighbors to vertices, or joined pairs of components along a single path. Thus the components at each stage were trees, and so the limit $T$ must be acyclic.

For each fixed $t$, the interesting case could only have occurred finitely many times as noted before the claim. Thus each vertex, once added to $C\inv(t)$, must have eventually experienced the uninteresting case, at which point it would be made degree $\Delta$ by construction. Since we added new vertices at the first stage for each $t$, and we always take these to be the least available, each $x \in X$ is eventually added to some $C\inv(t)$, and subsequently made to have degree $\Delta$. Thus $T$ is truly $\Delta$-regular and its vertex set is $X$, which was chosen to be computable. Similarly the edge set of of $T$ is $E$, also chosen to be computable.

Finally, $T$ is clearly computable: The construction of the $T_n$'s was clearly recursive, so to check whether a given $(x,e) \in X \times E$ is in $T$, we can run the construction until $x$ is added to the vertex set, at which point it gets $\Delta$ incident edges, and we can check if any of these is $e$.
\end{proof}

\begin{claim}
$T$ has no computable $\Pi$-coloring.
\end{claim}

\begin{proof}
Suppose $\phi_t$ is a $\Pi$-coloring of $T$, and consider our construction for this fixed value of $t$. In between each interesting case stage of our construction, we have finitely many $\Delta$-regular paths $P_i$, and are waiting for $\phi_t$ to converge to a $\Pi$-coloring of them. Since $\phi_t$ is in fact a $\Pi$-coloring of $T$, we are always guaranteed that this will eventually happen. That is, the interesting case will occur unboundedly many times for this $t$, but as noted before Claim \ref{claim:computable}, it can only occur $|\mathcal{V}| < \omega$ times.
\end{proof}

Thus $T$ is as desired.

\end{proof}

\end{document}
